\documentclass[11pt,a4paper]{article}
\usepackage[margin=25mm]{geometry}
\usepackage{amsmath,amssymb,amsthm}
\setlength{\emergencystretch}{2em}
\newtheorem{lemma}{Lemma}
\newtheorem{proposition}[lemma]{Proposition}
\newtheorem{corollary}[lemma]{Corollary}
\newcommand{\gen}[1]{\langle#1\rangle}
\title{Forests obtained by contracting an involution}
\author{Human Compression project: a component of the L044 recovery}
\date{Author version; manuscript review pending}
\begin{document}\maketitle

We isolate an elementary reduction for the two-block case in
Cameron--Fernandes--Leemans--Mixer, \emph{Highest rank of a polytope
for $A_n$}, arXiv:1605.09173v1, Proposition~3.5. Contraction below
produces a graph. It does not require every generator to permute the
contracted pairs. The final equality exclusions are separate obligations.

\paragraph{Hypotheses and notation.}
Let $I=\{0,\ldots,r-1\}\setminus\{i\}$, and let
$H=\gen{\rho_j:j\in I}$ be transitive on two preserved blocks
$B_0,B_1$, each of size $m$. Every $\rho_j$ is an even involution,
and $\rho_j\rho_k=\rho_k\rho_j$ for $|j-k|>1$.
Suppose $H_j=\gen{\rho_k:k\in I\setminus\{j\}}$ is
intransitive for every $j\in I$. Choose $\ell\in I$ such that
$l=\rho_\ell$ interchanges the blocks. Its orbits are $m$ pairs,
each meeting both blocks once. In particular $m$ is even. Put
\[
 C=\{j\in I\setminus\{\ell\}:\rho_jl=l\rho_j\},\qquad
 R=I\setminus(C\cup\{\ell\}).
\]
The string relations imply $|R|\le2$. For an involution, an actual
edge means one of its transpositions. Contract an actual edge to
the unordered pair of its $l$-orbits, deleting loops. Multiplicity
is retained only when explicitly stated. An $R$ generator need
not induce a permutation of these $m$ vertices.

For each $j\in I\setminus\{\ell\}$ choose an actual $j$-edge
joining different $H_j$-orbits; it exists by transitivity of $H$.
Every $H_j$-orbit is a union of $l$-pairs. Let $e_j$ be the
contracted edge so chosen, and let $F$ consist of these $r-2$ edges
on all $m$ vertices.

\begin{lemma}\label{forest}
The graph $F$ is a forest, with no parallel edges. If an actual
$h$-edge, $h\in I\setminus\{\ell\}$, joins two different
vertices in one component of $F$, those vertices are precisely
the endpoints of $e_h$.
\end{lemma}
\begin{proof}
Every walk in the contracted graph using labels other than $j$
lifts to an actual walk using those labels and $l$: within each
intermediate pair insert an $l$-edge if necessary to reach the
endpoint of the next actual edge. The endpoints can likewise be
adjusted to any prescribed points in their pairs. Since $l\in H_j$,
such a walk cannot join the endpoints of the chosen actual $j$-edge.
Any cycle of selected edges would give precisely such a walk by
removing its unique $j$-edge. This also excludes a cycle of two
parallel selected edges.

Now take the unique $F$-path between the endpoints of a contracted
$h$-edge. If that path contained $e_j$ with $j\ne h$, the new
$h$-edge and the rest of the path would bypass $e_j$ without using
$j$, a contradiction by the same lifting argument. The path is
therefore the single edge $e_h$, since only one edge of that label
was selected.
\end{proof}

\begin{proposition}\label{rankbound}
If $R\ne\varnothing$, then $r\le m$. At equality $F$ has exactly
two components.
\end{proposition}
\begin{proof}
A generator $h=\rho_j$ with $j\in R$ joins at least two distinct
unordered pairs of $l$-pairs. To see this, first suppose $h$ preserves
each block. There are no actual $h$-edges inside an $l$-pair. If
there were just one nonloop contracted edge, its multiplicity would
be one or two. Evenness forces two, and the resulting simultaneous
transpositions on the two rows commute with $l$.

If $h$ interchanges the blocks, it is a perfect matching. A single
nonloop contracted edge must be doubled: once one point of a pair
is matched outside that pair, its other point cannot be matched
internally. On all other pairs $h$ equals $l$. This again commutes
with $l$. Zero nonloop edges also imply commutation in both cases.

If $F$ were connected, Lemma~\ref{forest} would allow only the
single contracted edge $e_j$ for this generator, contrary to the
preceding paragraph. Thus $F$ has at least two components, whence
$r-2\le m-2$. Equality gives exactly two components.
\end{proof}

\begin{lemma}\label{local}
Suppose $r=m$ and $R\ne\varnothing$. For any
$h\in I\setminus\{\ell\}$, its simple contracted graph consists
of $e_h$ alone, $e_h$ and one edge across the two components of
$F$, or a triangle containing $e_h$ and one vertex in the other
component. For $h\in R$ there are precisely the following forms.
\begin{enumerate}
\item If $\rho_h$ preserves the blocks, it is a double
transposition. Its two contracted edges are distinct, one selected
and one crossing the component cut.
\item If $\rho_h$ interchanges the blocks, its contracted graph
is a triangle with multiplicity one on each edge. On every other
$l$-pair it equals $l$. On the three exceptional pairs, $l\rho_h$
acts by a three-cycle on one row and its inverse on the other.
\end{enumerate}
\end{lemma}
\begin{proof}
Write $e_h=ab$ in one component of $F$. All additional distinct
contracted $h$-edges cross the component cut by Lemma~\ref{forest}.
Given two such edges $cd,ef$, orient them with $c,e$ in the first
tree and $d,f$ in the second. The tree paths from $c$ to $e$ and
from $d$ to $f$, together with these two edges, form a cycle.
Either tree path is allowed to have length zero. Every selected
edge on these paths must have label $h$, by the lifting argument.
The path in the component not containing $ab$ therefore has
length zero, and the other path consists of $ab$. The two new
edges are consequently $ax,bx$ for one vertex $x$ in the other
tree. A third distinct crossing edge cannot satisfy this condition
with each of those two edges. This proves the first assertion.

For a block-preserving generator, the contracted edges are the
union of the matchings in the two rows. When the simple edges are
distinct this gives a proper colouring by two colours, so it
cannot contain an odd cycle. A generator in $R$ thus has exactly
two simple edges. Their multiplicities are at most two. Evenness
allows $(1,1)$ or $(2,2)$; the latter interchanges whole pairs and
commutes with $l$. Only $(1,1)$ remains. There can be no other
transpositions inside pairs, since these would interchange rows.

For a block-interchanging generator, each pair has either an
internal matching edge or two outgoing incidences. Its nonloop
contracted multigraph is therefore 2-regular on its nonisolated
vertices. With at most two simple edges it would consist of
doubled disjoint edges, giving commutation with $l$. Hence it is
the triangle, with every multiplicity one. Outside it the matching
is internal. Choose coordinates $(v,0),(v,1)$ with $l$ the row
interchange. An involution interchanging rows has the form
$(v,0)\mapsto(p(v),1)$, $(v,1)\mapsto(p^{-1}(v),0)$.
Here $p$ is a three-cycle on the exceptional vertices and fixes
all others. The asserted description of $l\rho_h$ follows.
\end{proof}

\begin{proposition}\label{bridges}
At the equality in Lemma~\ref{local}, let $J$ contain all simple
contracted edges of the $R$ generators, on all $m$ vertices,
including isolates. Let $E_C=\{e_j:j\in C\}$.
Every edge of $E_C$ is a bridge of $J\cup E_C$.
If $|R|=2$, then $\operatorname{rank}(J)=3$ and $J\cup E_C$
is connected. Here graph rank is the number of vertices minus
the number of components.
If $|R|=1$, then $J\cup E_C$ is also connected; it is a tree
when the sole $R$ generator preserves the two blocks.
\end{proposition}
\begin{proof}
Any alternate path around $e_j$, $j\in C$, in $J\cup E_C$
uses only $R$ labels and the other selected $C$ labels. It lifts
to a forbidden $H_j$ walk. The same argument excludes a duplicate
of $e_j$ among the edges of $J$. Adding the $|C|=m-2-|R|$
bridges thus raises graph rank by exactly $|C|$. Since the final
rank is at most $m-1$, we have
$\operatorname{rank}(J)\le |R|+1$.

Each individual $R$ graph has rank two by Lemma~\ref{local}.
For $|R|=1$, adding the $m-3$ bridges gives rank $m-1$, hence
connectedness. In the block-preserving case $J$ has just two
edges and is a forest; adjoining bridges gives a tree.

For $|R|=2$, suppose their union still had rank two. The component
partition of each individual graph must then equal that of the
union. Either both graphs consist of the same two disjoint edges,
or both are supported on the same three vertices. In the first
case exactly one edge is internal to an $F$-component, so both
labels must select that same edge. This contradicts
Lemma~\ref{forest}. In the second case the three vertices cannot
all lie in one $F$-component, because each $R$ graph has a crossing
edge. Their distribution is $2+1$, with exactly one internal
unordered pair. Both graphs must contain and select that pair,
giving the same contradiction. Thus $\operatorname{rank}(J)>2$.
The upper bound already proved gives rank three; adjoining the
$m-4$ bridges gives rank $m-1$ and connectedness.
\end{proof}

\paragraph{Attribution and scope.}
The chosen edges between omission orbits and the use of string
relations belong to the source argument. The present proof supplies
the explicit walk lifting, local support statements and bridge-rank
reduction, without assuming a quotient action on the pairs. These
steps were written by the root agent and checked in working form
by a distinct agent. No computational census is a premise.
This component deliberately assumes $R\ne\varnothing$ for its
rank bound. The separate empty-$R$ reduction requires the parent
alternating C-group, the acquired transitive rank theorem, finite
bases and the smaller-degree induction in the pair-block theorem;
preservation of pairs alone does not justify that invocation.
No whole-original recovery, open-conjecture solution, formal proof
certificate or numerical compression claim follows here.
\end{document}
